\section{Conclusiones}
\label{sec:conclusiones}

Este trabajo desarrolló un modelo de Programación Lineal Entera Mixta para
estudiar el número mínimo de aplicaciones activas de la capa no lineal
$\chi$ en versiones reducidas de Keccak.

La formulación representa bit a bit dos ejecuciones simultáneas de la
permutación y calcula sus diferencias mediante restricciones XOR exactas.
Las transformaciones $\theta$, $\rho$, $\pi$, $\chi$ e $\iota$ se modelaron
de acuerdo con la estructura de una ronda de Keccak, mientras que las
variables binarias $a_{r,y,k}$ permitieron identificar y minimizar la
actividad de cada aplicación local de $\chi$.

El modelo se aplicó a tamaños de palabra

\[
z\in\{4,8\},
\]

correspondientes a estados de 100 y 200 bits y a 20 y 40 aplicaciones
locales de $\chi$ por ronda, respectivamente. Los experimentos consideraron
una, dos y tres rondas.


\subsection{Respuesta al objetivo del trabajo}
\label{subsec:respuesta_objetivo}

El objetivo principal consistía en modelar Keccak reducido mediante MILP
y determinar o acotar el número mínimo de S-boxes activas para diferentes
tamaños de palabra y números de rondas.

La implementación desarrollada responde a este objetivo mediante:

\begin{itemize}
    \item una representación exacta de las transformaciones de Keccak;
    \item dos ejecuciones simultáneas para calcular diferencias XOR;
    \item variables binarias para identificar las S-boxes activas;
    \item una función objetivo que minimiza la actividad total;
    \item problemas auxiliares de factibilidad para establecer cotas;
    \item testigos concretos reconstruibles mediante una implementación
    funcional independiente.
\end{itemize}

La correspondencia entre intentos y rondas se utilizó únicamente como
criterio experimental. No se interpretó como una modificación interna de
Keccak ni como una relación criptográfica demostrada entre intentos de un
atacante y profundidad de la permutación.


\subsection{Resultados principales}
\label{subsec:conclusiones_resultados}

Para una ronda se certificó

\[
N_{\mathrm{act}}^{\min}(4,1)
=
N_{\mathrm{act}}^{\min}(8,1)
=
1.
\]

La cota inferior se deriva de la imposibilidad de que una diferencia no
nula desaparezca después de las transformaciones lineales. La existencia
de testigos con una sola S-box activa completa la prueba de optimalidad
global.

Para dos rondas se obtuvo

\[
N_{\mathrm{act}}^{\min}(4,2)
=
N_{\mathrm{act}}^{\min}(8,2)
=
4.
\]

CBC demostró que no existe una trayectoria con actividad total menor o
igual que tres, mientras que los testigos reconstruidos alcanzaron la
distribución

\[
2+2.
\]

La coincidencia entre ambas evidencias permite afirmar que cuatro es el
mínimo global para los dos tamaños de palabra.

Para tres rondas se establecieron los intervalos

\[
5
\leq
N_{\mathrm{act}}^{\min}(4,3)
\leq
13,
\]

y

\[
5
\leq
N_{\mathrm{act}}^{\min}(8,3)
\leq
14.
\]

Las cotas inferiores son globales y se derivan del mínimo exacto de dos
rondas junto con la necesidad de activar al menos una S-box adicional en
la tercera ronda.

Las cotas superiores están respaldadas por testigos con distribuciones

\[
2+2+9
\]

para $z=4$, y

\[
2+2+10
\]

para $z=8$.

Los valores 13 y 14 son óptimos dentro de la familia restringida
$2+2+c$, pero no constituyen mínimos globales certificados para el espacio
completo de tres rondas.


\subsection{Interpretación}
\label{subsec:conclusiones_interpretacion}

Los resultados muestran que una diferencia puede permanecer localizada
durante una ronda, pero la difusión acumulada incrementa la cantidad mínima
de componentes no lineales activos cuando se consideran más rondas.

La igualdad de los mínimos para $z=4$ y $z=8$ en una y dos rondas indica
que duplicar el tamaño de palabra y el número de aplicaciones locales de
$\chi$ no obliga, en estos horizontes, a aumentar la actividad mínima.

En tres rondas, los mejores testigos restringidos difieren en una S-box.
Sin embargo, como las cotas inferiores son iguales y los extremos superiores
provienen de una familia específica, no puede concluirse que el mínimo
global aumente necesariamente con el tamaño de palabra.

El número de S-boxes activas debe interpretarse como una medida estructural
de propagación no lineal. Esta cantidad no determina por sí sola la
probabilidad diferencial de una trayectoria ni la seguridad completa de
SHA-3.


\subsection{Aporte metodológico y computacional}
\label{subsec:aporte_metodologico}

La principal ventaja de la formulación pareada es que la propagación a
través de $\chi$ surge de dos evaluaciones exactas de la función booleana.
Esto evita depender exclusivamente de una aproximación diferencial
truncada y permite recuperar parejas concretas de estados.

CBC fue suficiente para:

\begin{itemize}
    \item obtener y validar testigos;
    \item demostrar infactibilidad en problemas auxiliares;
    \item certificar los mínimos globales de una y dos rondas;
    \item verificar estados y soportes obtenidos mediante enumeración.
\end{itemize}

Para tres rondas fue necesario complementar el modelo MILP con una búsqueda
exhaustiva restringida. Esta combinación permitió obtener testigos
verificables sin atribuirles un nivel de optimalidad superior al respaldado
por las evidencias.

La implementación se acompañó de una referencia funcional independiente,
artefactos estructurados y una suite de 245 pruebas automatizadas. Estos
elementos fortalecen la trazabilidad y permiten reconstruir los resultados
principales.


\subsection{Limitación principal}
\label{subsec:conclusion_limitacion}

La principal limitación es que el mínimo global de tres rondas permanece
abierto.

La búsqueda realizada es exhaustiva dentro de la familia

\[
2+2+c,
\]

pero no considera todas las distribuciones posibles de actividad durante
las tres rondas. En consecuencia, el resultado correcto para estos
experimentos es un intervalo de certificación y no un único valor exacto.

Esta distinción es esencial: los testigos de actividades 13 y 14 demuestran
que existen trayectorias con dichos valores, pero no descartan la existencia
de trayectorias globalmente mejores.


\subsection{Trabajo futuro}
\label{subsec:trabajo_futuro}

Las extensiones más relevantes del trabajo son:

\begin{enumerate}
    \item explorar distribuciones de tres rondas diferentes de
    $2+2+c$;

    \item incorporar restricciones de ruptura de simetrías y desigualdades
    válidas que fortalezcan la relajación del modelo;

    \item comparar CBC con solvers de mayor rendimiento, como Gurobi;

    \item desarrollar estrategias de descomposición que separen la
    generación de soportes, la asignación de valores absolutos y la
    validación MILP;

    \item incorporar pesos diferenciales de $\chi$ para minimizar el peso
    de una trayectoria y no únicamente el número de S-boxes activas;

    \item evaluar tamaños de palabra y números de rondas mayores.
\end{enumerate}


\subsection{Conclusión final}
\label{subsec:conclusion_final}

El número mínimo de S-boxes activas en Keccak reducido puede estudiarse de
manera sistemática mediante una combinación de modelado MILP, problemas de
factibilidad, enumeración controlada y validación independiente.

Para una y dos rondas se obtuvieron mínimos globales exactos. Para tres
rondas se construyeron cotas globales y testigos reproducibles, manteniendo
una distinción explícita entre optimalidad global y optimalidad restringida.

En conjunto, la implementación proporciona una respuesta rigurosa a la
actividad planteada y establece una base extensible para futuros estudios
de propagación diferencial en Keccak.


\section*{Disponibilidad del código y los artefactos}
\addcontentsline{toc}{section}{Disponibilidad del código y los artefactos}
\label{sec:disponibilidad_codigo}

El código fuente, las pruebas automatizadas, el notebook y los artefactos
experimentales se encuentran disponibles en:

\begin{center}
\url{https://github.com/hsancheza90-art/keccak-milp-criptografia}
\end{center}

La versión utilizada para los resultados del informe está identificada
mediante:

\begin{center}
\texttt{keccak-milp-experiments-v1}
\end{center}

y el commit:

\begin{center}
\texttt{0b089aa}.
\end{center}

Los comandos de reproducción, la descripción de los archivos generados y
el procedimiento de validación se presentan en la
Sección~\ref{sec:reproducibilidad}.